\section{Results}

\subsection{Zero-Shot Baseline Evaluation}

Table~\ref{tab:baselines} reports zero-shot accuracy across 11 VLMs.
Random chance is 12.5\% for 8-option questions and 25.0\% for role
identification (4 options).

\begin{table}[t]
\centering
\caption{Zero-shot accuracy (\%) on primary questions by category.
Best per column in \textbf{bold}.
$^\dagger$AWQ 4-bit quantized.}
\label{tab:baselines}
\small
\begin{tabular}{lcccccc}
\toprule
Model & Params & Simple & Compound & Complex & Ident. & Overall \\
\midrule
InternVL2.5-8B~\cite{chen2024internvl25}            & 8B  & --   & --   & --   & --   & --   \\
InternVL3-9B~\cite{zhu2025internvl3}                & 9B  & --   & --   & --   & --   & --   \\
InternVideo2.5-Chat-8B~\cite{wang2025internvideo25}  & 8B  & --   & --   & --   & --   & --   \\
Ovis2.5-9B~\cite{lu2024ovis}                        & 9B  & --   & --   & --   & --   & --   \\
Ovis2.5-9B-Thinking~\cite{lu2024ovis}               & 9B  & --   & --   & --   & --   & --   \\
Qwen2.5-VL-7B~\cite{bai2025qwen25vl}                & 7B  & --   & --   & --   & --   & --   \\
VideoLLaMA3-7B~\cite{zhang2025videollama3}           & 7B  & --   & --   & --   & --   & --   \\
LLaVA-Video-7B-Qwen2~\cite{zhang2024llavavideo}      & 7B  & --   & --   & --   & --   & --   \\
Qwen3-VL-8B-Instruct~\cite{bai2025qwen3vl}          & 8B  & --   & --   & --   & --   & --   \\
Qwen3-VL-8B-Thinking~\cite{bai2025qwen3vl}          & 8B  & --   & --   & --   & --   & --   \\
Qwen2-VL-72B-AWQ$^\dagger$~\cite{wang2024qwen2vl}   & 72B & --   & --   & --   & --   & --   \\
\bottomrule
\end{tabular}
\end{table}

% KEY ARGUMENTS to make in this paragraph (fill when numbers arrive):
% 1. Even the best model barely clears 46% — the task is genuinely hard.
% 2. Complex questions (aggressor+action+victim, sequence) drop below 34%
%    for ALL models — compositional reasoning is the bottleneck.
% 3. Thinking-mode variants (Ovis-Thinking, Qwen3-Thinking) show no
%    improvement over their base counterparts — extended inference-time
%    compute does not help without task-specific supervision.
% 4. The 72B model does not dominate the 7-9B models, suggesting the
%    challenge is representational, not capacity-limited.

\subsection{Where Models Fail: Compositional Reasoning}

Table~\ref{tab:per_type_primary} breaks down accuracy by question type
for the top baselines. The pattern is consistent: models identify
individual attributes (aggressor, victim, action) at 40--56\%, but
accuracy collapses when they must compose these attributes into a
coherent interaction.

\begin{table}[t]
\centering
\caption{Per-question-type zero-shot accuracy (\%) for the top-3 baselines,
grouped by compositional demand. Accuracy drops sharply as questions
require binding more attributes together.}
\label{tab:per_type_primary}
\small
\setlength{\tabcolsep}{4pt}
\begin{tabular}{llccc}
\toprule
Tier & Question Type & Model A & Model B & Model C \\
\midrule
\multirow{4}{*}{Single attribute}
  & Primary action          & -- & -- & -- \\
  & Aggressor identification & -- & -- & -- \\
  & Victim recognition       & -- & -- & -- \\
  & Role identification      & -- & -- & -- \\
\midrule
\multirow{2}{*}{Two attributes}
  & Action + victim          & -- & -- & -- \\
  & Aggressor + victim       & -- & -- & -- \\
\midrule
\multirow{2}{*}{Three+ attributes}
  & Aggressor + action + victim & -- & -- & -- \\
  & Sequence verification       & -- & -- & -- \\
\bottomrule
\end{tabular}
\end{table}

% KEY ARGUMENT: The accuracy drop from single → two → three attributes
% is monotonic across all models. This validates the tiered question
% design and shows the benchmark isolates compositional reasoning as a
% specific failure mode, not just general difficulty.

\subsection{Progressive Fine-Tuning Ablation}

Table~\ref{tab:ablation} shows the contribution of each training stage.

\begin{table}[t]
\centering
\caption{Stage-wise ablation on InternVL2.5-8B (20\% training data).
$\Delta$ is absolute improvement over zero-shot.}
\label{tab:ablation}
\small
\begin{tabular}{lccccc}
\toprule
Stage & Simple & Compound & Complex & Overall & $\Delta$ \\
\midrule
Zero-shot (base)   & -- & -- & -- & -- & ---   \\
+ SFT              & -- & -- & -- & -- & --    \\
+ CoT-SFT          & -- & -- & -- & -- & --    \\
+ ADPO             & -- & -- & -- & -- & --    \\
\bottomrule
\end{tabular}
\end{table}

% KEY ARGUMENTS (fill when numbers arrive):
% 1. SFT provides the largest single jump — format compliance matters.
% 2. CoT specifically improves compound/complex questions (where chains
%    are applied) while leaving simple questions flat — targeted gain.
% 3. ADPO adds a further boost on exactly the categories where the
%    hardest distractors live (role reversal, wrong action).
% 4. All three stages are needed; no single stage subsumes the others.

\subsection{Distractor Hardness Validation}

The hardness taxonomy (Section~\ref{sec:pairs}) predicts that role
reversals should fool models most often and cross-video distractors
least. Table~\ref{tab:hardness_errors} tests this prediction
empirically.

\begin{table}[t]
\centering
\caption{Distractor selection rate (\%) by hardness category.
Higher = the model is more frequently fooled. Categories ordered by
predicted difficulty.}
\label{tab:hardness_errors}
\small
\begin{tabular}{lcc}
\toprule
Distractor Category & Zero-shot & + ADPO \\
\midrule
Role reversal             & -- & -- \\
Wrong action              & -- & -- \\
Wrong victim              & -- & -- \\
Wrong aggressor           & -- & -- \\
Bystander substitution    & -- & -- \\
Cross-video               & -- & -- \\
\bottomrule
\end{tabular}
\end{table}




